\section{The residual rank estimate at the two distinguished primes}

% BEGIN J81H1
\paragraph{J81H1. A bound on all character blocks without a rank lifting hypothesis.}
Keep the native family hypotheses and the common hub of J80H2.
At a twisted index with $p=q$ or $d$ and a good point $\pi$,
let $C_\rho,C_0,\ldots,C_{t-1}$ be the global cocycle modules
of I81G1. Write $h_E$ for the class number and $u_e$ for the
native length of the $e$th positive twist.
The selected point satisfies the killing and finiteness guards of
J78H2. Put $Q_0=f^2t^2(r+t)$ and $u=\log b$.
With $e_O=\operatorname{length}_O(O/p)$, one has
\[
 \frac{1}{e_O}\sum_{0\le e<t}\operatorname{length}_O C_e
 \ \le\ \sum_{0<e<t}u_e+v_p(h_E)+p^{-2}+180Q_0u .
\]
Indeed K81C6 and H80M4 give all of $\mathrm{TBAt}$ at the joint
$c_B$ of H80M3, for every required good coefficient extension.
The character calculation in I81G5--I81G6 used no rank row estimate.
Its error is at most $t c_B\sum_w(d_w+1)$. By R79Q4,
$\sum_w d_w=2f$ and the number of places is at most $f$.
Insert $c_B\le60tf(r+t)u$. The auxiliary-prime term costs at
most $p^{-2}$ by the same strict degree guard and killing condition
as in I81G6.

In particular let $D_\rho$ denote the image of the simultaneous
restriction of $C_\rho$ in
the direct sum over $w\mid p$ of the local quotients by their
finite parts. Put $\delta_\rho=\operatorname{length}_O D_\rho/e_O$.
D81S2 shows this is a finite rational number for each point.
The kernel of this restriction injects in the Selmer group by
I81G6's kernel argument, which also holds without a row bound.
Consequently
\[
 v_{\rm CC}(i,\pi)\le s_\pi+\sum_{0<e<t}u_e+v_p(h_E)
                 +p^{-2}+180Q_0u+\delta_\rho.
\]
There is only one joint image here. D81S2 also gives a
point-dependent upper bound for it in terms of elementary divisors
of the global action on the common column module.
% END J81H1

% BEGIN J81H2
\paragraph{J81H2. The ramified-prime field on the original family hub.}
Let $S=S_c$ and $B=b/p+S+t+\eta+u$ as in H80M5. To obtain the
ramified field on the hub it suffices, at every guarded point of
every sufficiently late member, to prove
\[
                 \delta_\rho\le2080 f(rB+mb)+f .
\]
For this reduction, apply J81H1 and use $u>1$ and $Q_0\ge f$. The polynomial
$181 f^2t^2(r+t)$ times $u$ absorbs $180Q_0u+f$. The polynomial
has coefficient sum $362$ and total degree $5$. Apply the comparison
of J78H2 at $A=2080$ with the same constant already chosen for J80H2.
The argument tests only the point and representation required by the
original ramified field.

U81L2 proves the pointwise row bound of H80M2 with constant
$160$ for every row. In particular its bound holds for each of
the maps of I81G3. I81G4 and subadditivity over the places therefore
give $\delta_\rho\le\sum_w B_w$. H80M5's summation yields
$\sum_w B_w\le13\cdot160 f(rB+mb)+f$, as required. This
proves the displayed bound on the same chosen family hub. It gives
the native $\mathrm{Concl.ram}$ there together with the twelve
fields of J80H2. The common constant can be retained, increasing
the late-member cutoff to include the finitely many parameter
inequalities used above. The two remaining fields in the denominator
family conclusion are $\mathrm{mid}$ and $\mathrm{two}$. The complete
global denominator contract additionally asks for the original family
construction and the final bound on its $\eta$-sum.
% END J81H2
